\begin{lemma}
\label{lemma-curve-affine-projective}
Let $X$ be a curve over $k$. Then either $X$ is an affine scheme or $X$
is H-projective over $k$.
\end{lemma}

\begin{proof}
Choose $X \to \overline{X}$ with
$\overline{X} \setminus X = \{x_1, \ldots, x_r\}$ as in
Lemma \ref{lemma-reduced-dim-1-projective-completion}.
Then $\overline{X}$ is a curve as well.
If $r = 0$, then $X = \overline{X}$ is H-projective over $k$.
Thus we may assume $r \geq 1$ and our goal is to show that
$X$ is affine. By Lemma \ref{lemma-open-in-affine-curve-affine}
it suffices to show that $\overline{X} \setminus \{x_1\}$ is affine.
This reduces us to the claim stated in the next paragraph.

\medskip\noindent
Let $X$ be an H-projective curve over $k$. Let $x \in X$ be a closed point
such that $\mathcal{O}_{X, x}$ is a discrete valuation ring. Claim:
$U = X \setminus \{x\}$ is affine. By Lemma \ref{lemma-regular-point-on-curve}
the point $x$ defines an effective Cartier divisor of $X$.
For $n \geq 1$ denote $nx = x + \ldots + x$ the $n$-fold sum, see
Divisors, Definition \ref{divisors-definition-sum-effective-Cartier-divisors}.
Denote $\mathcal{O}_{nx}$ the structure sheaf of $nx$ viewed as a
coherent module on $X$.
Since every invertible module on the local scheme $nx$ is trivial
the first short exact sequence of
Divisors, Remark \ref{divisors-remark-ses-regular-section}
reads
$$
0 \to \mathcal{O}_X
\xrightarrow{1} \mathcal{O}_X(nx) \to \mathcal{O}_{nx} \to 0
$$
in our case. Note that $\dim_k H^0(X, \mathcal{O}_{nx}) \geq n$.
Namely, by Lemma \ref{lemma-chi-tensor-finite} we have
$H^0(X, \mathcal{O}_{nx}) = \mathcal{O}_{X, x}/(\pi^n)$ where $\pi$ in
$\mathcal{O}_{X, x}$ is a uniformizer and the powers $\pi^i$ map to
$k$-linearly independent elements in $\mathcal{O}_{X, x}/(\pi^n)$
for $i = 0, 1, \ldots, n - 1$. We have $\dim_k H^1(X, \mathcal{O}_X) < \infty$
by Cohomology of Schemes, Lemma
\ref{coherent-lemma-proper-over-affine-cohomology-finite}.
If $n > \dim_k H^1(X, \mathcal{O}_X)$
we conclude from the long exact cohomology sequence
that there exists an $s \in \Gamma(X, \mathcal{O}_X(nx))$
which is not a section of $\mathcal{O}_X$.
If we take $n$ minimal with this property, then $s$ will
map to a generator of the stalk
$\left(\mathcal{O}_X(nx)\right)_x$ since otherwise it would
define a section of $\mathcal{O}_X((n - 1)x) \subset \mathcal{O}_X(nx)$.
For this $n$ we conclude that $s_0 = 1$ and $s_1 = s$
generate the invertible module $\mathcal{L} = \mathcal{O}_X(nx)$.

\medskip\noindent
Consider the corresponding morphism
$f = \varphi_{\mathcal{L}, (s_0, s_1)} : X \to \mathbf{P}^1_k$
of Constructions, Section \ref{constructions-section-projective-space}.
Observe that the inverse image of $D_{+}(T_0)$ is $U = X \setminus \{x\}$
as the section $s_0$ of $\mathcal{L}$ only vanishes at $x$.
In particular, $f$ is non-constant, i.e., $\Im(f)$ has more
than one point. Hence $f$ must map the generic
point $\eta$ of $X$ to the generic point of $\mathbf{P}^1_k$.
Hence if $y \in \mathbf{P}^1_k$ is a closed point, then $f^{-1}(\{y\})$
is a closed set of $X$ not containing $\eta$, hence finite.
Finally, $f$ is proper\footnote{Namely, a H-projective variety
is a proper variety by Morphisms, Lemma
\ref{morphisms-lemma-projective-is-quasi-projective-proper}.
A morphism of varieties whose source is a proper variety is a proper morphism
by Morphisms, Lemma \ref{morphisms-lemma-image-proper-scheme-closed}.}.
By Cohomology of Schemes, Lemma
\ref{coherent-lemma-proper-finite-fibre-finite-in-neighbourhood}\footnote{One
can avoid using this lemma which relies on the theorem of formal
functions. Namely, $X$ is projective hence it suffices to show
a proper morphism $f : X \to Y$ with finite fibres between quasi-projective
schemes over $k$ is finite. To do this, one chooses an affine open of $X$
containing the fibre of $f$ over a point $y$ using that any finite set of
points of a quasi-projective scheme over $k$ is contained in an affine.
Shrinking $Y$ to a small affine neighbourhood of $y$ one reduces to the
case of a proper morphism between affines. Such a morphism is finite by
Morphisms, Lemma \ref{morphisms-lemma-integral-universally-closed}.}
we conclude that $f$ is finite. Hence $U = f^{-1}(D_{+}(T_0))$
is affine.
\end{proof}